% eden-compressed.tex
% The Four Streams
% Intended as an include fragment under master.tex.
% LuaLaTeX; relies on the Hebrew/font/source primitives defined in welibc.tex.

\Apocryphon{The Four Streams}

\begin{center}
{\large The Nameless\par}
\vspace{0.35em}
{\small Genesis 2:11--14, resegmentation, and the possibility of a deliberately multiplexed text\par}
\end{center}

\section*{Abstract}

Genesis 2:10--14 describes one river leaving Eden and dividing into four heads: Pishon, Gihon, Tigris, and Euphrates. The received reading is coherent and almost certainly primary. This paper asks a different question: whether part of the same consonantal sequence was also composed to support secondary word division.

Three hypotheses are tested. \textbf{A} begins a secondary stream at the \heb{שם} near the end of 2:11. \textbf{B} begins at the \heb{שם} in 2:12. \textbf{C} treats both as active, making the two locative instances of \heb{שם} latent name-like restart markers between the explicit first and second river headers. Together with the received text, these are the four streams of the title.

The strongest evidence is not a hidden sentence. It is a constrained, zero-edit resegmentation: the consonants stay in order, none are added or removed, and shifted boundaries repeatedly produce legitimate Hebrew or Northwest Semitic packets. The most striking core is the surface \heb{טוב שם הבדלח ואבן}, which admits \heb{טוב שמה בדל חוא בן}; a shorter restart yields \heb{שם הבדל חוא בן}. Here \heb{שמה} is richly ambiguous, \heb{בדל} is an attested noun ``severed piece,'' \heb{הבדל} is an attested Hiphil form of \heb{בדל}, and \heb{בן} is ordinary ``son, offspring.'' The comparative interest of \heb{חוא} is primarily Phoenician \textit{ḥwʾ}, ``live,'' not a proposed spelling of Eve.

The case is cumulative. Genesis 2:10 has already said that the river ``divides'' (\heb{יפרד}); \heb{שם} repeatedly means both ``name'' and ``there''; the fourth river alone lacks a \heb{שם} header; and A reproduces the main lexical anchors of the Pishon description in altered order. Nearby wordplay and orthographic variation make such attention to letters plausible, but do not prove it. The result is therefore graded: B is the strongest local resegmentation, A the strongest long carry-chain, and C the most explanatory literary model. Whether the combined pattern exceeds chance is a computational question that remains open.

\section*{1.\quad The four streams}

The relevant surface text is ordinary Hebrew. Pishon surrounds Havilah, where gold is found; the gold is good; \heb{בדלח} and \heb{שהם} are also there. Their exact material referents are uncertain, but the syntax is not.\cite{BDB,Westermann1984}

\Table{p{0.25\linewidth}p{0.65\linewidth}}{%
\textbf{Hebrew} & \textbf{English} \\ \hline
\heb{ומשם יפרד} & and from there it divides \\
\heb{והיה לארבעה ראשים} & and it became four heads \\
\heb{שם האחד פישון} & the name of the first is Pishon \\
\heb{אשר שם הזהב} & where the gold is \\
\heb{וזהב הארץ ההוא טוב} & and the gold of that land is good \\
\heb{שם הבדלח ואבן השהם} & there are bdellium and the shoham stone \\
\heb{ושם הנהר השני גיחון} & and the name of the second river is Gihon \\
\heb{ושם הנהר השלישי חדקל} & and the name of the third river is Tigris \\
\heb{והנהר הרביעי הוא פרת} & and the fourth river is Euphrates
}

Nothing here requires emendation. The proposal concerns \emph{resegmentation} or \emph{rebracketing}: a secondary parsing superimposed on a secure primary one.

\begin{description}
\item[\textbf{A.}] Restart at the late 2:11 \heb{שם}, in \heb{אשר שם הזהב}.
\item[\textbf{B.}] Restart at the 2:12 \heb{שם}, in \heb{טוב שם הבדלח}.
\item[\textbf{C.}] Treat both restarts as active.
\end{description}

I use \emph{multiplexing} as a computer-science analogy for this arrangement. The claim is not that the text encodes a second continuous plaintext. It is that one ordered consonantal stream may support more than one meaningful tokenization.

\section*{2.\quad Controls}

Resegmentation is easy to overfit in unpointed Hebrew. Short roots, one-letter prefixes, productive suffixes, and homography create many accidental words. The useful question is therefore not whether another word can be found, but how constrained the alternate reading is.

The analysis uses six rules:

\begin{enumerate}
\item consonants remain in transmitted order;
\item none are inserted, deleted, or rearranged;
\item final letter forms are normalized only graphically, not phonologically;
\item no Paleo-Hebrew look-alike substitutions are used in the primary argument;
\item Biblical Hebrew forms outrank comparative Semitic rescue forms;
\item lexical validity, syntax, local relevance, and analyst freedom are scored separately.
\end{enumerate}

``Zero-edit'' below means exactly this: after ignoring spaces, vocalization, cantillation, and the positional shapes of final letters, the consonantal sequence is unchanged.

Nor should the proposal be defended by claiming that early Hebrew lacked word division. Northwest Semitic inscriptions attest deliberate separation of words; universal early \textit{scriptio continua} is not a safe assumption.\cite{Millard1970,Tov2012} If the effect is ancient, it crosses a primary boundary deliberately. It is literary resegmentation, not recovery of a forgotten spacing.

\section*{3.\quad The surface architecture}

The river list establishes \heb{שם} as a naming marker, then omits it only before the fourth river.

\Table{lll}{%
\textbf{River} & \textbf{Header} & \textbf{Name} \\ \hline
First & \heb{שם האחד} & Pishon \\
Second & \heb{ושם הנהר השני} & Gihon \\
Third & \heb{ושם הנהר השלישי} & Tigris \\
Fourth & \heb{והנהר הרביעי} & Euphrates
}

The omission is harmless as ordinary prose. Under C, however, it becomes relevant because two additional \heb{שם} tokens have already appeared inside the expanded Pishon description.

The Pishon and Gihon clauses also preserve a small orthographic asymmetry:

\Table{p{0.30\linewidth}p{0.60\linewidth}}{%
\textbf{Hebrew} & \textbf{English} \\ \hline
\heb{הוא הסבב את כל ארץ החוילה} & it is the one surrounding all the land of Havilah \\
\heb{הוא הסובב את כל ארץ כוש} & it is the one surrounding all the land of Cush
}

The same participial function is defective in 2:11 and plene in 2:13. Plene/defective spelling varies throughout the Bible, and how much of the Masoretic distribution reaches back to first composition remains disputed.\cite{Elitzur2023,Hornkohl2024} The pair is therefore not proof that an author placed every mater deliberately. It is simply a local fact: two nearly parallel clauses use different consonantal spellings for the same grammatical role. Any reading sensitive to matres should notice it.

\section*{4.\quad Stream A}

A begins with the late 2:11 \heb{שם}. Its main interest is mechanical.

\Table{p{0.43\linewidth}p{0.47\linewidth}}{%
\textbf{Surface} & \textbf{A} \\ \hline
\heb{שם הזהב} & \heb{שם הזה} + \heb{ב} \\
\heb{ב + וזהב} & \heb{בוזה} + \heb{ב} \\
\heb{ב + הארץ} & \heb{בה} + \heb{ארץ} \\
\heb{ארץ + ההוא} & \heb{ארצה + הוא} \\
\heb{שם + הבדלח} & \heb{שמה + בדל + ח} \\
\heb{ח + ואבן} & \heb{חוא + בן}
}

This is a \emph{carry-chain}: a consonant released by one displaced boundary becomes material for the next token. That dependency is what distinguishes A from a bag of unrelated substrings.

The full candidate sequence is:

\Table{p{0.16\linewidth}p{0.25\linewidth}p{0.49\linewidth}}{%
\textbf{Token} & \textbf{Type} & \textbf{Defensible value} \\ \hline
\heb{שם} & noun / adverb & name / there \\
\heb{הזה} & demonstrative & this, this one \\
\heb{בוזה} & participle & despising, treating as contemptible \\
\heb{בה} & prep. + suffix & in/on/by/with her or it \\
\heb{ארצה} & directional noun / verb & toward land; consonantally also I will favor \\
\heb{הוא} & pronoun & he, it \\
\heb{טוב} & adjective & good, favorable \\
\heb{שמה} & adverb / noun / verb & there; her name; she placed; desolation \\
\heb{בדל} & noun / root & severed piece; separate/distinguish \\
\heb{חוא} & comparative packet & Phoenician life-associated form \\
\heb{בן} & noun & son, offspring \\
\heb{השה} & noun & the flock animal
}

The weakness is syntax. \heb{בוזה} and \heb{בה} are individually good forms, but \heb{בוזה בה} is not a natural Biblical Hebrew way to say ``despises her''; \heb{בזה} normally takes a direct object or other constructions.\cite{BDB} A should therefore not be presented as a fluent hidden sentence.

Its opening is still structurally interesting. The catalogue has just introduced Pishon with \heb{שם האחד}; a few words later the same consonants recur in \heb{אשר שם הזהב}. On the surface the second \heb{שם} means ``there.'' At a secondary level it can still be heard as ``name,'' giving low-commitment openings such as \heb{שם הזה}, ``the name of this [one] \ldots,'' or retaining \heb{שם הזהב} as a name/gold packet. Neither is a complete clause; both fit the naming register already established.

The strongest part of A begins later:

\Table{p{0.43\linewidth}p{0.47\linewidth}}{%
\textbf{Surface} & \textbf{Alternative} \\ \hline
\heb{טוב} & \heb{טוב} \\
\heb{שם} & \heb{שמה} \\
\heb{הבדלח} & \heb{בדל} + \heb{ח} \\
\heb{ואבן} & \heb{חוא} + \heb{בן}
}

Several useful facts align here. The article \heb{ה} moves left and produces \heb{שמה}; the remainder \heb{בדל} is an attested Biblical Hebrew noun, ``severed piece,'' in Amos 3:12; the released \heb{ח} combines with \heb{וא} to give \heb{חוא}; and the residue is ordinary \heb{בן}, ``son, offspring.''\cite{BDB}

The comparative packet \heb{חוא} should carry limited weight. Biblical Hebrew writes Eve \heb{חוה}. The relevant old comparison is Phoenician \textit{ḥwʾ}, ``live,'' which BDB places beside the Hebrew life complex.\cite{BDB} This makes the packet interesting near Genesis 3's life language without requiring an alternative spelling of Eve.

A also reproduces the main lexical anchors of the Pishon description in altered order:

\Table{p{0.30\linewidth}p{0.60\linewidth}}{%
\textbf{Surface frame} & \heb{שם} $\rightarrow$ \heb{הוא} $\rightarrow$ \heb{ארץ} $\rightarrow$ \heb{שם} \\
\textbf{A frame} & \heb{שם} $\rightarrow$ \heb{ארצה} $\rightarrow$ \heb{הוא} $\rightarrow$ \heb{שמה}
}

The resemblance is structural, not syntactic identity. The outer \heb{שם}/\heb{שמה} frame remains while the two central anchors reverse.

\section*{5.\quad Stream B}

B starts later and therefore asks less of the text. Two exact segmentations are available.

\Table{p{0.18\linewidth}p{0.30\linewidth}p{0.42\linewidth}}{%
\textbf{Stream} & \textbf{Segmentation} & \textbf{Best-supported values} \\ \hline
B1 & \heb{שם הבדל חוא בן} & name/there; separation-form; life-associated packet; son \\
B2 & \heb{שמה בדל חוא בן} & there/name/placed; severed piece; life-associated packet; son
}

B1 retains the \heb{ה} with \heb{בדל}. The exact form \heb{הבדל} is attested as a Hiphil infinitive absolute of \heb{בדל}; Isaiah 56:3 uses it in the semantic field of separation.\cite{BDB,JouonMuraoka2006} B2 shifts the \heb{ה} left, producing the richly ambiguous \heb{שמה} and the exact noun \heb{בדל}, ``severed piece.''

Neither should be forced into a polished proposition. Their strongest shared kernel is simply:

\begin{center}
\heb{שם}/\heb{שמה} \quad + \quad separation.
\end{center}

That matters because 2:10 has just introduced the whole passage with:

\Table{p{0.25\linewidth}p{0.65\linewidth}}{%
\textbf{Hebrew} & \textbf{English} \\ \hline
\heb{ומשם} & and from there \\
\heb{יפרד} & it divides \\
\heb{והיה לארבעה ראשים} & and it became four heads
}

The roots \heb{פרד} and \heb{בדל} are different, but the semantic field is the same and the correspondence is immediate. B is therefore the strongest local hypothesis: it begins at a salient \heb{שם}, requires few boundary decisions, avoids A's weak opening syntax, and lands directly on the idea that defines the paragraph.

\section*{6.\quad Stream C}

C treats A and B as two restarts generated by the same local mechanism.

\Table{p{0.13\linewidth}p{0.36\linewidth}p{0.40\linewidth}}{%
\textbf{Position} & \textbf{Surface function} & \textbf{Secondary possibility} \\ \hline
1 & \heb{שם האחד פישון} & explicit first name header \\
A & \heb{אשר שם הזהב} & latent \heb{שם הזה\ldots} / name-gold restart \\
B & \heb{טוב שם הבדלח} & latent \heb{שם הבדל\ldots} or \heb{שמה בדל\ldots} \\
2 & \heb{ושם הנהר השני גיחון} & explicit second name header \\
3 & \heb{ושם הנהר השלישי חדקל} & explicit third name header \\
4 & \heb{והנהר הרביעי הוא פרת} & fourth river; no \heb{שם}
}

The surface text already alternates two values of the same packet: \heb{שם} means ``name'' in the river headers and ``there'' inside the Pishon excursus. C proposes that the alternation is functional. The locative uses can also act as latent restart markers.

This does not produce six rivers. It produces a denser header structure inside the description of the first. The fourth river's missing \heb{שם} then becomes conditionally relevant: not evidence by itself, but exactly the sort of asymmetry one might expect if the passage has already spent the delimiter twice inside the Pishon section.

In computational language, C treats \heb{שם} as an \emph{overloaded delimiter}: name marker in one parse, locative in another, possible restart marker in a third. C has the greatest explanatory reach, and therefore also the greatest danger of post-selection.

\section*{7.\quad Establishing the protocol}

The wider case is that Genesis 2--3 may teach the reader, very early, which kinds of textual relation are worth noticing. ``Protocol'' is deliberately imported terminology: not a claim about ancient cryptography, but a model of how a text can establish decoding conventions before reusing them.

Several local features fit that model.

\paragraph{Nested consonantal forms.}
The opening contains the family \heb{אד}, \heb{אדם}, \heb{אדמה}. They do not occur as a simple chronological derivation---\heb{אדם} and \heb{אדמה} precede \heb{אד}---but they form an obvious nested string family, and the narrative explicitly exploits \heb{אדם}/\heb{אדמה}.

\paragraph{Overt paronomasia.}
The story openly plays with \heb{איש}/\heb{אשה}, then moves from \heb{ערומים}, ``naked,'' in 2:25 to \heb{ערום}, ``crafty,'' in 3:1; Genesis 3:20 associates \heb{חוה} with life.\cite{WaltkeOConnor1990,Westermann1984} This does not prove hidden word division. It establishes that lexical resemblance is already part of the composition's visible technique.

\paragraph{Name, place, placement.}
The river catalogue sits between the placement of the human in the garden and the naming of the animals. The consonants \heb{שם} participate in the nearby semantic fields of ``name,'' ``there,'' and the placement root. \heb{שמה} expands the ambiguity further. A high-ambiguity packet is repeatedly placed at high-salience positions.

\paragraph{Matres.}
The local \heb{הסבב}/\heb{הסובב} pair shows that near-equivalent forms need not be spelled identically in the received consonantal text. Because transmission complicates authorship of individual matres, this is supporting context, not primary evidence.

\paragraph{Gihon.}
Locally, \heb{גיחון} resembles the snake's \heb{גחון}, ``belly,'' in Genesis 3:14. Remotely, Gihon is the site of Solomon's anointing in 1 Kings 1. Under Friedman's reconstruction of an extended early prose work running into the court history, these belong near opposite ends of one large composition.\cite{Friedman1998} That reconstruction is debated, but if granted provisionally, Gihon becomes a plausible local/remote lexical bridge.

The clearest remote echo is the death warning:

\Table{p{0.27\linewidth}p{0.29\linewidth}p{0.34\linewidth}}{%
\textbf{Text} & \textbf{Hebrew} & \textbf{Structure} \\ \hline
Genesis 2:17 & \heb{ביום אכלך ממנו מות תמות} & on the day you do X $\rightarrow$ dying you shall die \\
1 Kings 2:37 & \heb{ביום צאתך \ldots ידע תדע כי מות תמות} & on the day you do X $\rightarrow$ knowing you shall know $\rightarrow$ dying you shall die
}

Friedman draws attention to this recurrence and to the surrounding knowledge, good/bad, and death language in the Shimei episode.\cite{Friedman1998} It does not prove A--C. It shows what long-range reuse of an early lexical configuration could look like.

Taken together, these features suggest a simple protocol: attend to small consonantal changes, repeated packets, shifted senses, spelling variation, and echoes that may be local or remote. Under that model the river section is a natural training ground. The story says one river becomes four; the text may be showing that one consonantal stream can also branch.

\section*{8.\quad Evidence, chance, and many weak signals}

The four streams can be ranked without pretending they are equally strong.

\Table{p{0.12\linewidth}p{0.21\linewidth}p{0.20\linewidth}p{0.20\linewidth}p{0.17\linewidth}}{%
\textbf{Stream} & \textbf{Mechanical constraint} & \textbf{Lexical quality} & \textbf{Syntax} & \textbf{Best case} \\ \hline
Surface & excellent & excellent & excellent & primary reading \\
A & long zero-edit carry-chain & mostly strong & uneven & strongest mechanism \\
B & short, few choices & very strong core & compressed & strongest local case \\
C & inherits A+B & inherits A+B & need not form one sentence & strongest literary model
}

No single observation settles intentionality. The case is a bundle of weak and medium signals: repeated \heb{שם}, attested \heb{בדל}/\heb{הבדל}, the carry-chain, the division vocabulary of 2:10, the river-header asymmetry, \heb{הסבב}/\heb{הסובב}, Gihon, and the Genesis--Kings death formula. Their force lies in partial alignment.

A physics analogy is constructive interference: many small contributions can reinforce one another without any one path being decisive.\cite{FeynmanHibbs1965} A closer modern analogy is a high-dimensional embedding, where a large dot product can arise from many weakly aligned dimensions rather than one explicit rule. The problem for literary analysis is to determine whether the alignment is in the text or in the reader.

That problem cannot be solved merely by invoking a favorite formalism. Neyman--Pearson testing requires specified sampling distributions.\cite{NeymanPearson1933} Bayesian inference, including Jaynes's extended-logic interpretation, requires meaningful likelihoods.\cite{Jaynes2003} Minimum Description Length is attractive because C gains value if one compact ``protocol'' compresses many observations, but MDL still requires a defensible description language.\cite{Rissanen1978} In each case the hard part comes first: translating a qualitative literary claim into a model.

That is the \textbf{compilation problem}. ``Same root'' is easy to formalize. ``Same syntactic construction'' is usually manageable. ``This passage appears to teach the reader how later correspondences should be noticed'' is not. Formalizing only what is easy to count risks measuring the residue and discarding the phenomenon.

A useful test of A--C should therefore freeze the rules before searching and compare Genesis 2 against several baselines: real sliding windows from the Hebrew Bible, shuffled controls, $n$-gram surrogates, and morphologically plausible synthetic text. Candidate segmentations should be scored separately for consonantal coverage, boundary coherence, lexical prior, morphosyntax, semantic locality, structural recurrence, and analyst intervention. Work on biblical parallel detection and source partition already shows that philological judgments can be subjected to computational comparison.\cite{NaaijerRoorda2016,YoffeEtAl2023}

The Transformer changed what can be attempted here.\cite{VaswaniEtAl2017} The point is not to ask a language model whether a pun was intended. That answer is not evidence. The change is representational: semantic similarity, long-range dependency, alternative tokenization, and style-preserving negative controls can now be modeled without reducing everything to a handful of hand-written features.

A serious future system could combine a character-level segmentation lattice, diachronically controlled morphology, syntax scores, contextual embeddings, explicit penalties for cross-language rescue, and held-out passages. It should not output ``J intended C: 97.3\%.'' It should answer narrower questions: How rare is a chain like this? Does it remain exceptional after linguistic controls? Does a score trained on independent Hebrew wordplay rank this passage highly without being tuned to it?

Those are empirical questions. They are now much less inaccessible than they were before 2017.

\section*{9.\quad Conclusion}

The surface reading of Genesis 2:11--14 is secure. The question is whether it is the only reading engineered into the consonants.

The evidence is strongest in three places. First, A gives a real zero-edit carry-chain across several displaced boundaries. Second, B produces an unusually compact \heb{שם}/\heb{שמה} + separation sequence immediately after the river has been said to divide. Third, C explains why the expanded Pishon description contains two extra \heb{שם} packets between explicit river-name headers and why the fourth river can arrive without one.

The case should not be inflated. A is not good continuous Hebrew. \heb{חוא} is supporting comparative evidence, not a key. The missing \heb{שם} before Euphrates is suggestive only under a larger model. Orthographic variation may reflect transmission as well as composition. And the total pattern has not yet been tested against a proper null corpus.

Nor should the case be diluted. The resegmentation is exact, locally concentrated, mechanically linked, and semantically entangled with the immediate passage. It deserves a stronger explanation than ``Hebrew has short words.'' The right next step is not to add looser substitutions until a secret sentence appears. It is to freeze the present zero-edit phenomenon and measure how exceptional it is.

If that kind of work succeeds, the implications extend beyond this passage. The Hebrew Bible has been read as scripture, law, history, liturgy, source-critical composite, and literary anthology. It may also contain works engineered at a density that our traditional categories do not describe well: narrative, sound, spelling, word boundary, local structure, and long-range recurrence acting at once.

That possibility should be tested against other ancient corpora, not protected from comparison. If the Bible proves unusually dense under fair tests, then a modern tradition centered on it may eventually look unlike any existing sect. Such a tradition would need the old attentiveness to every letter without the old requirement to know in advance why every letter matters. It would keep the text central while making interpretation revisable; it would treat philology, archaeology, source criticism, statistics, and machine learning as ordinary tools of reading; and it would regard uncertainty not as disloyalty but as metadata.

Its basic discipline would be simple: distinguish the text from the model, the primary reading from the secondary possibility, evidence from dogma, and reproducible pattern from private revelation.

The first river divides into four.

The first task is to learn how many streams we are actually reading.

\begin{thebibliography}{99}

\bibitem{BDB}
Francis Brown, S. R. Driver, and Charles A. Briggs.
\newblock \textit{A Hebrew and English Lexicon of the Old Testament}.
\newblock Oxford: Clarendon Press, 1907.

\bibitem{Elitzur2023}
Joel Elitzur.
\newblock ``Plene Spelling and Defective Spelling in the Hebrew Bible: The Question of Dating.''
\newblock \textit{Journal of the American Oriental Society} 143, no. 4 (2023): 745--765.
\newblock doi:10.7817/jaos.143.4.2023.ar027.

\bibitem{FeynmanHibbs1965}
Richard P. Feynman and Albert R. Hibbs.
\newblock \textit{Quantum Mechanics and Path Integrals}.
\newblock New York: McGraw--Hill, 1965.

\bibitem{Friedman1998}
Richard Elliott Friedman.
\newblock \textit{The Hidden Book in the Bible}.
\newblock San Francisco: HarperSanFrancisco, 1998.

\bibitem{Hornkohl2024}
Aaron D. Hornkohl.
\newblock ``Orthography.''
\newblock In \textit{Diachronic Diversity in Classical Biblical Hebrew}.
\newblock Cambridge: Open Book Publishers, 2024.

\bibitem{Jaynes2003}
E. T. Jaynes.
\newblock \textit{Probability Theory: The Logic of Science}.
\newblock Edited by G. Larry Bretthorst.
\newblock Cambridge: Cambridge University Press, 2003.

\bibitem{JouonMuraoka2006}
Paul Joüon and Takamitsu Muraoka.
\newblock \textit{A Grammar of Biblical Hebrew}.
\newblock 2nd ed. Rome: Pontifical Biblical Institute, 2006.

\bibitem{Millard1970}
A. R. Millard.
\newblock ```Scriptio Continua' in Early Hebrew: Ancient Practice or Modern Surmise?''
\newblock \textit{Journal of Semitic Studies} 15, no. 1 (1970): 2--15.
\newblock doi:10.1093/jss/15.1.2.

\bibitem{NaaijerRoorda2016}
Martijn Naaijer and Dirk Roorda.
\newblock ``Parallel Texts in the Hebrew Bible: New Methods and Visualizations.''
\newblock 2016.
\newblock arXiv:1603.01541.

\bibitem{NeymanPearson1933}
Jerzy Neyman and Egon S. Pearson.
\newblock ``On the Problem of the Most Efficient Tests of Statistical Hypotheses.''
\newblock \textit{Philosophical Transactions of the Royal Society of London, Series A} 231 (1933): 289--337.
\newblock doi:10.1098/rsta.1933.0009.

\bibitem{Rissanen1978}
Jorma Rissanen.
\newblock ``Modeling by Shortest Data Description.''
\newblock \textit{Automatica} 14, no. 5 (1978): 465--471.
\newblock doi:10.1016/0005-1098(78)90005-5.

\bibitem{Tov2012}
Emanuel Tov.
\newblock \textit{Textual Criticism of the Hebrew Bible}.
\newblock 3rd ed. Minneapolis: Fortress Press, 2012.

\bibitem{VaswaniEtAl2017}
Ashish Vaswani, Noam Shazeer, Niki Parmar, Jakob Uszkoreit, Llion Jones,
Aidan N. Gomez, Łukasz Kaiser, and Illia Polosukhin.
\newblock ``Attention Is All You Need.''
\newblock In \textit{Advances in Neural Information Processing Systems 30}, 2017.
\newblock arXiv:1706.03762.

\bibitem{WaltkeOConnor1990}
Bruce K. Waltke and M. O'Connor.
\newblock \textit{An Introduction to Biblical Hebrew Syntax}.
\newblock Winona Lake, IN: Eisenbrauns, 1990.

\bibitem{Westermann1984}
Claus Westermann.
\newblock \textit{Genesis 1--11: A Commentary}.
\newblock Translated by John J. Scullion.
\newblock Minneapolis: Augsburg, 1984.

\bibitem{YoffeEtAl2023}
Gideon Yoffe, Axel Bühler, Nachum Dershowitz, Israel Finkelstein,
Eli Piasetzky, Thomas Römer, and Barak Sober.
\newblock ``A Statistical Exploration of Text Partition Into Constituents:
The Case of the Priestly Source in the Books of Genesis and Exodus.''
\newblock 2023.
\newblock arXiv:2305.02170.

\end{thebibliography}
